% Original-Hume contribution to planned164.2; distinct from Lange 2011.
\chapter{Hume and the Passage from Experience to Expectation}
\label{sec:p15-hume-and-the-problem-of-induction}

Hume's question concerns a familiar achievement: from something observed,
we believe something not observed. A letter supports a belief about an
absent friend; a manufactured object suggests former inhabitants; bread
previously eaten encourages an expectation about another piece. The
inferred fact need not be in the future. Historical testimony, a present
but distant event and an anticipated effect all extend beyond the immediate
deliverances of sense and memory.

The central argument in Sections~IV--V of the \emph{Enquiry concerning
Human Understanding} asks what explains and supports this passage.
Hume distinguishes demonstrative relations, experienced conjunctions,
the expectation of an unobserved conjunction, and the customary movement
of thought producing that expectation. His negative argument and his
positive account answer different questions. Failure to discover a
noncircular argument is not the discovery that people can stop relying
on experience, nor a mathematical theorem that every inductive procedure
fails \parencite[E~4.3--4.5, 4.14--4.23, 5.1--5.9]{Hume1777EnquiryHTO}.

The presentation here follows the \emph{Enquiry}'s final 1777 copytext.
Its first publication in 1748, the earlier \emph{Treatise}, and later
philosophers' reconstructions are different presentations. Paragraph
locators beginning with E identify the \emph{Enquiry}. Sections~II--III
supply the idea and association distinctions used by the account;
VI--VII clarify probability and necessary connexion; IX and
XII.20--XII.34 qualify the animal-reasoning and sceptical conclusions.
Hume's illustrations and psychological assertions are historical
philosophical evidence, not newly measured populations or experiments.
The earlier history, from Aristotle's \emph{epagoge} to Leibniz, is
reconstructed in \cref{sec:p15-w02b-p15-026-milton}.

\section{Two kinds of enquiry and the scope of the question}
\label{sec:p15-hume-fork}

\begin{hbdefinition}{def:p15-hume-kinds}{Hume's relations of ideas and matters of fact}
In E4.1, \emph{relations of ideas} are the objects of intuitive or
demonstrative certainty, exemplified by geometry, algebra and arithmetic.
Their truth does not require the existence of corresponding physical
objects. \emph{Matters of fact} concern existence and occurrences.
Their contraries can be conceived without the contradiction which would
defeat a demonstratively false assertion. Evidence for a matter of fact
can be very strong without being of the former kind.

Hume's later term \emph{moral reasoning} in this division means reasoning
about fact and existence. It does not mean only ethical argument.
His term \emph{proof} can also mean an experiential argument leaving no
ordinary room for doubt, as distinct from a demonstration. The note to
Section~VI explicitly distinguishes demonstrations, experiential proofs
and less certain probabilities.
\end{hbdefinition}

The sunrise example does not say that tomorrow's sunrise and its failure
are equally likely. It says that denying the former does not involve
the contradiction found in denying a demonstrated relation. Conceivability
is used to exclude one form of necessity, not to assign equal numerical
weights to every conceived alternative
\parencite[E~4.1--4.2, 6.0n10]{Hume1777EnquiryHTO}.

The question begins where direct observation ends. A current sound can
support a belief about its source, a watch on an island about its maker,
and heat about a fire's other effects. Hume calls the connecting relation
causal, allowing direct, remote and collateral links. In particular,
inferring one collateral effect from another is not automatically the
claim that the first effect produces the second. An inferential connection
and a direct production relation must not be conflated.

There are consequently at least three things to assess. What was observed?
What further fact is believed? What connexion is invoked between them?
Simply repeating the observation as the conclusion does not explain
the extension; simply naming a connexion does not establish it.
The formal explanation relation of
\cref{sec:p16-finite-abduction} is a different, explicitly defined task:
there the background and admissible assumptions are supplied. Hume is
asking about our reliance on the empirical connexions themselves.

\section{Why the effect is not read off from the cause}
\label{sec:p15-hume-experience}

Hume first denies that inspection of an entirely new object's sensible
qualities reveals its causes or effects. His imagined Adam possesses
strong rational faculties but lacks experience. Transparency and fluidity
do not tell that novice what water will do; brightness and warmth do not
tell him every consequence of fire. These are philosophical contrast
cases separating possession of concepts from acquaintance with an
empirical operation, not observations of an actual first human
\parencite[E~4.6--4.8]{Hume1777EnquiryHTO}.

Unfamiliar cases make the dependence on experience easier to notice.
Hume mentions smooth marble surfaces with differing resistance to
separation, gunpowder, magnets and nourishment in differently constituted
animals. Familiar cases can conceal that dependence. Someone used to
billiards may think the second ball's movement is already contained in
the thought of the first ball's impact. Hume replies that the two
movements are distinct occurrences. Imagining contact alone does not
exclude rest, rebound or another direction. The example concerns what
is supplied by the description considered independently of experience;
it does not deny consequences under an already accepted physical law.

The same distinction explains his treatment of mathematics in natural
philosophy. Geometrical reasoning can calculate dimensions or consequences
once an empirical mechanical principle is supplied. Its precision does
not, by itself, discover or justify that principle. Hume's historical
mechanical terminology is not adopted as a current equation identifying
force with momentum. His point is the conditional role of mathematics:
deduction applies the assumed law; the law's empirical application needs
another basis \parencite[E~4.9--4.13]{Hume1777EnquiryHTO}.

Reducing several effects to fewer general principles can increase
understanding without revealing an ultimate productive power. Hume
mentions elasticity, gravity, cohesion and impulse as the kinds of
principles to which enquiry may arrive. A successful reduction therefore
must not be mistaken for the removal of every question about why those
principles hold. This is an epistemic limit within his account, not a
claim that mathematical modelling or further experiments are useless.

\section{The two steps in the demand for a foundation}
\label{sec:p15-hume-uniformity}

Section~IV's second part asks a further question. Even granting that
knowledge of causal relations comes from experience, why does that
experience support a new instance? The earlier answer has moved the
issue rather than settled it: causal reasoning appeals to experience,
but the passage from experienced to unexperienced cases still needs
explanation \parencite[E~4.14--4.16]{Hume1777EnquiryHTO}.

\begin{hbexample}{def:p15-hume-bread}{The observed bread and the next piece}
Hume separates the report that bread with certain sensible qualities
nourished someone on the occasions observed from the expectation that
another apparently similar piece will nourish that person later.
The new claim refers to another object or occasion. Repeating the
first report, however many observations it summarizes, does not turn
it into the second claim. Apparent similarity alone does not establish
unchanged hidden constitution or unchanged effects.

The example is historical task material. It supplies no nutritional
measurement, clinical recommendation or guarantee about actual food.
Its role is to expose the additional expectation, not to advise an
experiment with a harmful substance.
\end{hbexample}

For a \emph{demonstrative} justification, the difficulty is that a change
in the course of nature is not contradictory. Hume illustrates this
with apparently snowy material having an unfamiliar taste or sensation,
and with trees flourishing in what had previously been their dormant
season. The observed past can be acknowledged while the new outcome
differs. Thus the relevant predictive connexion is not established by
abstract demonstration alone. This objection does not say that a
deduction from additional premises about the laws and circumstances
cannot be valid. It asks where the additional empirical premises obtain
their authority \parencite[E~4.18]{Hume1777EnquiryHTO}.

For a \emph{probable} or experiential justification, Hume's objection
depends on his preceding account of factual inference. Such an argument
uses a causal connexion learned from experienced conjunctions and extends
that connexion to a new case. If it is offered to vindicate the very
extension it already employs, it presupposes the disputed expectation.
Past success of the practice is still a fact about the past. Taking it
as assurance of future success uses another instance of the same passage.
That is the circularity at issue, not the mere occurrence of experience
as a premise \parencite[E~4.19--4.21]{Hume1777EnquiryHTO}.

\begin{hbwarning}{def:p15-hume-question-boundary}{Which conclusion follows from the dilemma?}
The demonstrative objection and the experiential objection have different
premises. The first excludes a necessary conclusion from the past report
alone. The second says that the proposed experiential vindication already
uses the projective connexion at issue. The second is not proved merely
by showing that an inductive conclusion is not deductively entailed.

The target is a purported argumentative foundation of the extension, under
Hume's division and account of factual reasoning. No theorem that every
possible justification is either a deduction or a circular argument is
introduced as an axiom here. A response challenging his classification,
his causal account, or the demand for this kind of foundation must give
its own argument. Naming such a response does not already answer him.
\end{hbwarning}

Hume's wording also resists a simple slogan about universal irrationality.
E4.16 allows, for the argument, that the new proposition may justly be
inferred; the immediate challenge asks for the reasoning which effects
the transition. E4.21 distinguishes the agent, satisfied in practice,
from the philosopher seeking its foundation. E4.22 acknowledges that
failing to find an argument does not, by itself, establish that none
exists. These qualifications leave interpretation work to be done.
The text should not be rewritten either as a bare deductive
nonentailment observation or as an unqualified instruction never to
believe anything unobserved.

\section{Replies examined within the argument}
\label{sec:p15-hume-replies}

One reply points to a larger number of uniform observations. Hume asks
what intermediate reasoning makes their number decisive. A demonstrated
property of a circle is not changed by inspecting many circles, whereas
expectations concerning events do change with repeated experience.
This is an argument for a difference in the source of assurance. It is
not the assertion that several observations contain no more information
than one, or that sample size cannot matter in a specified statistical
model \parencite[E~4.20, 5.5]{Hume1777EnquiryHTO}.

Another reply says experience reveals the hidden powers responsible for
the sensible qualities. Hume replies that this restates the question
unless it supplies a new connexion. The old observations concern those
objects at those times; relating the familiar appearance of another
object to the same hidden powers is precisely the further inference.
Explaining a repeated regularity by a proposed mechanism can be
scientifically useful, but it does not follow that the mechanism's
continued applicability is demonstrated merely by naming it.

A practical reply points to successful conduct. Hume does not deny that
we act on experience. He distinguishes that fact from an explanation
of the authority of the passage. Asserting that the practice continues
therefore does not refute the philosophical question; asserting that
the question remains does not refute the practice's occurrence.

Finally, infants and animals learn to expect consequences without
possessing an elaborate philosophical proof. Hume argues that an
intricate argument proposed as the actual source of their expectation
would fail that explanatory task. This bears on the \emph{origin} of
the transition. A later reflective argument might assess a belief
without having caused its initial formation. Nor does the illustration
rule out every unreflective computation or every modern cognitive
mechanism. Hume provides no controlled task, recruited sample or
process-identification experiment establishing such a universal
exclusion \parencite[E~4.22--4.23, 9.5]{Hume1777EnquiryHTO}.

The negative argument thus has both an argumentative and a psychological
aspect. These must not be collapsed. A description of how an expectation
arises is not automatically its justification; a proposed justification
is not automatically a faithful description of the process producing it.

\section{Custom, belief and the preservation of practice}
\label{sec:p15-hume-custom}

Section~V first warns that philosophy can rationalize an existing
temperament rather than correct it: an ideal of self-sufficiency can
foster selfishness, and reflection on worldly vanity can excuse
indolence. Hume contrasts academic doubt, caution and restricted
enquiry with that use of doctrine. This is his practical and
normative framing, not an experimental finding that one school
always produces virtue or another always produces vice.

Hume's positive answer begins with an experienced person who has observed
many similar conjunctions. On encountering one member again, that person
expects its usual attendant. The change is a propensity acquired through
repetition, not newly perceived access to the objects' hidden productive
power. Hume calls this principle \emph{custom} or \emph{habit}. He
expressly says the name identifies an acknowledged kind of operation
without giving its ultimate cause
\parencite[E~5.2--5.6]{Hume1777EnquiryHTO}.

The principle supplies an explanation of why repeated experience affects
expectation and why the agent continues to infer even after being
convinced by the philosophical difficulty. It also explains a practical
dependence: without expectations extending beyond present impressions,
we could not ordinarily select means for ends or anticipate consequences.
That dependence is not itself a demonstration that the next expectation
will be true.

The expectation must start from something present to sense or memory.
A ruin can support an inference about earlier inhabitants; reading
historical testimony can lead through a chain of reports to alleged
observations. A merely hypothetical chain with no factual starting point
does not, for Hume, establish the actual existence of its objects.
The hypothesis and the observation supporting its application are
different parts of the account.

\begin{hbdefinition}{def:p15-hume-belief}{Belief as a manner of conceiving}
In E5.10--5.13, belief differs from an imagined fiction not by an
additional idea which can be annexed voluntarily, but by a characteristic
feeling or manner of conception. Hume describes the believed object as
conceived more firmly, steadily and forcefully. The imagination can
combine familiar materials into a fictional creature without thereby
making its existence believed.

This is Hume's phenomenological and psychological explication. Its
vivacity terms do not define a numerical probability scale, a calibrated
confidence measure or a criterion of truth.
\end{hbdefinition}

His comparisons bring together a relation and a present impression.
A friend's picture can enliven the friend's idea; approaching home can
strengthen its familiar associations; an object once used by a deceased
person or that person's child can revive memories. Religious images
and relics are further examples in his discussion. The force of his
religious polemic is not adopted as an independent verdict about their
users. The relevant proposed operation is the passage from a present
object to a related, more forceful conception.

These illustrations do not acquire their objects' existence simply
by enlivening an idea. Hume himself notes that the picture comparison
presupposes belief in the friend's existence. He then seeks an
analogous explanation of that factual belief. The firewood and sword
examples apply customary expectation to anticipated effects, not to
an arbitrary connection established only by imagination
\parencite[E~5.14--5.20]{Hume1777EnquiryHTO}.

E5.21--5.22 presents the operation as fitted to ordinary human life and
less dependent on slow reflective reasoning than such a basic capacity
could be. This is his explanatory proposal, not a proof from usefulness
to truth. Section~VI considers competing effects, and E12.22 expressly
allows that an instinct can be deceptive. Those qualifications prevent
turning the natural mechanism into a guarantee that every acquired
association is reliable.

The long note to E5.5 adds another qualification. Ordinary discourse
contrasts reason with experience, yet ostensibly reflective reasoning
in politics or practical life ultimately uses experienced general
principles too. Reflection can identify relevant circumstances and
trace consequences, while experience supplies and improves the maxims.
Hence his explanation does not exclude all attentive thought from an
experienced person's reasoning. It denies that such thought has a
wholly experience-independent empirical starting point.

\section{Probability and differences between reasoners}
\label{sec:p15-hume-probability}

Section~VI distinguishes the probability associated with considered
chances from the probability associated with variable observed effects.
In the die illustration, several equally possible sides bear one mark
and fewer bear another. More views converge on the former event, and
Hume proposes a correspondingly stronger expectation. The equal-side
description is part of the illustration; an actual die's distribution
is not determined by counting marks without that assumption.

In variable causal experience, the imagination must consider the
different effects actually encountered. Hume contrasts an apparently
uniform history with irregular medicinal effects and climate-dependent
weather. These are historical illustrations, not present treatment or
weather guidance. Conceiving an event as having occurred frequently
can strengthen its expectation without making competing occurrences
logically impossible. His statement that chance reflects ignorance of
causes is part of his philosophical account, not a contemporary theorem
excluding objective stochastic models
\parencite[E~6.1--6.4, 6.0n10]{Hume1777EnquiryHTO}.

The transfer of experienced proportions is consequently another use
of the passage under examination. Section~VI is not a derivation of
frequency convergence from bare observations. A proved law of large
numbers requires the law and sampling assumptions specified in its
own theorem. Calling an experiential assurance a proof also does not
convert it into a demonstration in the earlier division.

Section~IX extends the explanatory account by analogy to animal learning.
Hume describes acquired avoidance, trained responses and different
conduct by experienced and inexperienced animals. He argues that their
ordinary expectation cannot be generated by the elaborate argument
which philosophers themselves struggle to find. This is an analogical
natural-history argument. It does not provide experimental discrimination
among association, representation, decision or inference mechanisms.
Its appeal to analogy is compatible with using ordinary experience to
investigate how we infer; it would not, by itself, be a noncircular
vindication of all analogy \parencite[E~9.1--9.6]{Hume1777EnquiryHTO}.

The note to E9.5 prevents the stronger misunderstanding that all
experienced reasoners must perform alike. It names nine sources of
difference: attention, memory and observation; comprehension of complex
causal situations; length of consequence chains; maintenance of clear
ideas; separation of relevant from extraneous circumstances; accuracy
in forming general maxims; experience and availability of analogies;
prejudice, education, passion and party; and the enlargement of experience
through testimony, books and conversation. These are the source's
proposed differences, not measured independent effect sizes.

The same note explains why one carefully conducted experiment can matter
to an already experienced investigator. A general habitual expectation
has already been acquired. The investigator is not the wholly
unexperienced newcomer of Section~V. This reconciles the two examples
without supplying a demonstration that the prior general habit is
correct or that one trial establishes every application.

\section{Necessary connexion and the limits of a causal definition}
\label{sec:p15-hume-connexion}

Section~VII asks where the idea of productive necessity comes from.
Its opening contrasts the precision of mathematical terms with the
obscurity of many reflective ideas. Mathematics can require longer
chains and wider comparisons, while short moral or metaphysical
arguments can fail through ambiguous terms. Natural enquiry also
depends on finding suitable observations. This motivates a careful
explication of power; it does not establish a general ranking of
the sciences or make a clear definition an empirical discovery.
Its method invokes the impression--idea distinction of Section~II:
an idea is ordinarily analysed through the more lively sensory or
reflective experience from which it is copied. This is a philosophical
proposal, not a definition proving the existence of a suitable
impression for every term. Hume's own missing-shade-of-blue example
qualifies an exceptionless simple-copy claim. Section~III likewise
offers resemblance, contiguity and causation as association principles
while acknowledging the difficulty of proving the list exhaustive.
Neither premise is silently converted into an axiom of this handbook
\parencite[E~2.3--2.9, 3.1--3.3, 7.1--7.5]{Hume1777EnquiryHTO}.

In that prerequisite argument, impressions include felt sensations,
passions and acts of willing; ideas include their less lively
recollections and combinations. The golden mountain combines familiar
materials, and an enlarged conception of wisdom supplies Hume's
example of an apparently more remote idea. He offers analysis into
simple components and the effects of absent sensory experience as
two supports for the copying proposal. These are his philosophical
observations, not a universal clinical finding about every person's
conceptual capacities. The missing shade is acknowledged as a
counterinstance to an exceptionless simple-copy rule. His note on
innateness distinguishes being original from being present at birth
or merely being natural. Changing those meanings would change the
argument rather than refute the same stated proposition.

The association examples likewise distinguish a picture leading to
its original, an apartment leading to adjacent rooms, and a wound
leading to anticipated pain. The note treats contrast as possibly
combining resemblance and causation through an object's annihilation
and its former existence. Hume marks completeness of the enumeration
as difficult to establish. These proposed psychological associations
are not entailment rules: a thought's bringing another thought to
mind does not prove the latter's truth.

External observation of one collision supplies succession, not a
perceived necessity compelling that particular effect. Reflection on
volition does not give Hume the missing perception either. The limits
of voluntary bodily control are learned through experience; feeling
that one commands a movement is not knowledge of the intervening
mechanism. His examples of impaired movement and his dated anatomical
language are philosophical illustrations, not clinical diagnoses or
present physiology. The control of ideas and passions is also limited
and variable. Observing that an idea follows a decision to think does
not reveal the productive power by which it was raised.

The effort felt when overcoming resistance supplies a further candidate
impression. Hume allows that it contributes to ordinary loose talk of
power, while denying that it explains every attributed power or predicts
its effects independently of experience. Power is attributed to
inanimate objects and to operations without the same felt resistance.
Even where effort occurs, what follows remains something learned
\parencite[E~7.6--7.20, 7.15n13]{Hume1777EnquiryHTO}.

An occasionalist reply assigns the true productive role to divine
volition and makes ordinary objects occasions rather than causes.
Hume gives two objections. First, the analogies used in ordinary
experience are being extended beyond their experiential sphere.
Second, replacing an obscure bodily power with an equally obscure
mental or divine power does not provide the absent impression or
explanation. His further providential comparison between delegated
and continuously exercised power is a separate philosophical argument,
not empirical evidence of either mechanism. The note's accounts of
Newton and the Cartesians remain Hume's historical attributions, not
complete treatments of their writings
\parencite[E~7.21--7.25, 7.25n16]{Hume1777EnquiryHTO}.

The positive proposal locates a new impression in the \emph{observer}.
Repeated similar conjunctions establish a customary anticipation.
The observed instances may be alike, yet the experienced transition
of thought differs between the novice and the accustomed observer.
Hume identifies that felt determination as the source of the idea
of necessary connexion. This accounts for an idea and an expectation;
it is not direct inspection of a hidden tie in the objects.

E7.29 supplies two explications of cause: a followed object together
with regular succession of similar objects, and a followed object
whose appearance habitually carries thought to its attendant.
The first is accompanied by a counterfactual gloss about absence of
the first object. Hume presents these as ways of expressing the
connexion available to us. They are not a proof that regularity,
customary expectation and every modern counterfactual definition are
extensionally identical. Indeed his earlier heat/light example allows
inference between collateral effects. The finite observationally equal
models in \cref{thm:p16-causal-nonidentification} make a different,
explicit mathematical point: observational agreement alone need not
determine an equation-replacement intervention law. That result's
domains and structural assumptions belong to its own presentation,
not to an alleged hidden formal theorem in Hume.

The vibrating-string/sound illustration applies the two explications
to a familiar causal statement. The further note distinguishes a
measure of effects from comprehension of productive power itself.
Using the words force or energy does not settle that comprehension;
transferring an internal sensation of effort to an external object
can instead conceal the question
\parencite[E~7.26--7.30, 7.29n17]{Hume1777EnquiryHTO}.

\section{Sceptical inquiry without the abolition of ordinary life}
\label{sec:p15-hume-scepticism}

Section~XII distinguishes popular sceptical complaints from the
philosophical difficulty. Human fallibility, disagreement and changes
of judgment are not, by themselves, sufficient to abolish the evidence
on which life proceeds. The deeper challenge repeats the absence of
an argument connecting experienced conjunctions with further instances.
It also distinguishes the operation of an instinct from its inability
to deceive \parencite[E~12.21--12.22]{Hume1777EnquiryHTO}.

Hume's objection to excessive Pyrrhonism is practical and purposive.
Ordinary action brings the agent back to expectation; a universal,
durable suspension of this kind would not sustain life. That is not
a deductive refutation of every sceptical objection. It explains why
the philosophical difficulty can survive while the supposed permanent
behavioural conclusion does not. The theory's tendency to produce
useful, lasting effects is another consideration, distinct from
the truth of a particular belief.

His mitigated scepticism therefore has two forms. It recommends
modesty, caution and attention to counterarguments rather than
precipitate dogmatism. It also limits enquiry to questions fitted
to our capacities and experience. Philosophy can continue as a
methodized and corrected reflection on ordinary life without
claiming a view beyond all experiential limitations
\parencite[E~12.23--12.26]{Hume1777EnquiryHTO}.

The concluding classification confines demonstration to quantity
and number and directs other factual enquiries to experience.
History and the sciences of general facts have different objects;
morals and criticism additionally concern sentiment and taste.
His theological remarks and closing library test belong to that
methodological and polemical conclusion. They are not this handbook's
admission rule for legitimate enquiry, a proof eliminating modern
logic, or an instruction to destroy books
\parencite[E~12.27--12.34, 12.29n35]{Hume1777EnquiryHTO}.

\begin{hbwarning}{def:p15-hume-normative-boundary}{Three questions which habit does not collapse}
The \emph{causal question} asks how an expectation is produced.
The \emph{epistemic question} asks what supports it as a belief
about an unobserved fact. The \emph{practical question} asks
what a person should do under particular goals, alternatives and
stakes. Hume's natural explanation, his foundational challenge
and his recommendations of caution contribute differently.

A widespread or unavoidable mechanism does not, by that fact
alone, establish truth or an ethical duty. Conversely, inability
to supply a wholly independent argumentative foundation does not
make every contextual comparison between better and worse
reasoning disappear. Section~IX's differences and Section~XII's
qualified recommendations explicitly retain such comparisons.
\end{hbwarning}

\section{Using the problem alongside formal inference methods}
\label{sec:p15-hume-reference-use}

The Bayesian procedure in \cref{sec:p15-bayes-reference-procedure}
starts from a prior, an observation mechanism and a future law.
Its valid calculation does not itself supply or vindicate those
empirical inputs. The probability identity and the normative
claim that one must update by that identity are also distinguished
in \cref{sec:p15-rat41-dynamics}. Thus adding numerical probabilities
does not, by itself, remove Hume's demand for the connexion.
It can make a conditional conclusion and its assumptions more precise.

Exchangeability and the mixture representation in
\cref{sec:p15-exchangeability-and-definetti} similarly concern a
specified joint law on the stated domain. A representation theorem
does not establish that a physical sequence satisfies the required
symmetry. The learning method of \cref{def:p15-occam-method} fixes
its class and representation before sampling and states the
fresh-input distribution to which its error conclusion applies.
Its proved guarantee does not assert success under every changed
distribution. Neither theorem is invalidated by noticing that
its assumptions need justification; neither is an unconditional
philosophical justification of them.

Structural causal modelling makes further empirical commitments
explicit: assignments, disturbance laws and invariance under a
chosen intervention. Its observational and intervention questions
are not simply interchangeable uses of the word cause. Hume's
distinctions help identify the place at which an observed
regularity, a selected model, a philosophical explication and
a defended manipulation assumption enter the inference.

\begin{hbmethod}{def:p15-hume-question-record}{Locate the extension and the claimed support}
First identify the observation or remembered report and the
unobserved conclusion. State whether the extension concerns a
new time, object, location, testimony chain or intervention.
Separate apparent similarity from the mechanism or stability
assumption invoked.

Then ask what work each proposed premise does. Is it an
experiential regularity, a formal consequence under a model,
a probable assessment, an explanation of the reasoner's
transition, or a normative recommendation? If past success
is offered to vindicate the extension, expose the further
expectation linking that success to the new occasion.
If a mathematical guarantee is offered, retain its stated
distribution, domain and observation assumptions.

Finally distinguish the conclusion obtained from the question
left open. A success case identifies the extra connexion
rather than concealing it in a renamed principle. Failure
cases infer a demonstration from repetition alone, infer
truth from the strength of habit, infer an action norm from
the most frequent outcome, or treat failure to find an
argument as proof that every possible answer is impossible.
\end{hbmethod}

The resulting use is neither an instruction to abandon
empirical enquiry nor a promise to settle the foundational
problem by an identity. It is a way to keep a conditional
method, its empirical application and its philosophical
interpretation answerable to different questions.

\index{Hume, David!induction}
\index{induction!experience and custom}
\index{custom!causal expectation}
\index{necessary connexion!Hume's explication}
\index{scepticism!mitigated and excessive}
